\section{Collateral-Backed NFTs}
\label{sec:collateral-backed}

\subsection{Staked Collateral}

As users collateralize their Animal NFT or Egg, they are locking in value into the NFT. The NFT's intrinsic value is based on scarcity, and collateral is multiplied by the APY of the animal for additional rewards. Users can resell the NFT for a higher value than the purchase price plus collateral plus rewards earned. Alternatively, users can burn the NFT and get back the collateral value plus the rewards. However, the only way to get back staked collateral and rewards earned is to burn the NFT or sell it on the Zoo Marketplace.

Zoo is designed to easily add collateral or stake any currency on any of the Zoo Labs NFTs by ``feeding,'' ``breeding,'' or by minting an older NFT version of the animal.

\subsection{Releasing Collateral}

Once a user stakes collateral into their NFT, that value is locked in. The collateral earns APY based on the NFT's rarity, which determines its rewards yield percentage. As a user earns rewards, the rewards become part of the collateral, which increases the NFT's reward. In order to liquidate collateral and its rewards, the user has two options:

\begin{enumerate}[label=\Alph*.]
\item \textbf{Sell and transfer the NFT} for the value of the NFT plus its intrinsic value. When the user sells their NFT on the marketplace, they get back its staked value and have the opportunity to make more based on its intrinsic value, which is determined by rarity within the ecosystem. If an NFT is scarce, the ROI will be higher versus an NFT that is widely available.

\item \textbf{Burn the NFT} directly, returning the staked value plus rewards earned to date. When the NFT is burned, the allotment to mint older animals (if possible) is gone forever, as the NFT itself is gone forever, removing one of its species from the deflationary ecosystem. A 20\% tax is imposed upon burning: 10\% goes to the \DAO{} and 10\% goes back to the liquidity pool.
\end{enumerate}

\subsection{Rewards}

Every NFT is a representation of an endangered species. The ecosystem represents the animal kingdom of endangered species, over 18{,}000. The amount of NFTs released and put into circulation is dependent on how rare and endangered these animals are in the real world. The first drop will consist of 1{,}800 hatchable eggs with 3 varieties of separate rarities. These NFTs earn rewards and extra boosts from collateral staked. When hatched, each egg reveals an animal which earns rewards for a year.

\subsection{Feeding (Adding Collateral)}

First, a user incubates an egg which hatches into an Origin Baby Animal. Similar to Tamagotchi, users can ``care'' for their pets by feeding them and see their animals age up. The ``feeding'' action is defined by adding any currency to the NFT and locking it in until the user decides to burn the NFT or sell it to another user. With any Animal, users can continue to feed it and eventually mint an older version by feeding it the same value as its base cost.

There are 7 generations of NFTs in the Zoo ecosystem. It starts with an egg, a baby, then a teen, and finally an adult. If a user collects two adults they can breed them to mint a new baby animal. Of each Origin Egg, users can mint 1{,}500+ animals, collect them all, or sell them on the Zoo Marketplace. Users can also breed the animals if they collect the same species of adults that still have breeds left. The Origin Animals can breed 7 times and each generation after will have 1 less breed possible. The ultimate objective of the game is to collect the highest yielding and most rare NFT within the Zoo ecosystem.

\subsection{Boosts}

Boosts are determined by collateral added, providing an additional reward on top of the daily reward that comes with the animal:

\begin{equation}
\text{Total Reward} = \text{Daily Reward} + (\text{Collateral} \times \text{Boost} \times \text{Days}/365)
\end{equation}

\noindent The boost tiers are:
\begin{center}
\begin{tabular}{ll}
\toprule
\textbf{Collateral Range} & \textbf{Boost} \\
\midrule
\$10--\$99   & 1.1$\times$ \\
\$100--\$999  & 1.2$\times$ \\
\$1{,}000--\$9{,}999 & 1.3$\times$ \\
\$10{,}000+  & 1.4$\times$ \\
\bottomrule
\end{tabular}
\end{center}

\subsection{Breeding}

NFT Animals can be bred with other NFT Animals once they reach their adult form, achieved by ``feeding'' the NFT Animal a specified amount of crypto. Each NFT Animal can be bred up to 7 times, and any two adult animals of the same species can be bred with one another to produce an NFT Egg.

Users must breed two adult versions of the same species of NFT Animal in order to successfully produce offspring, which comes in the form of a totally random NFT Egg that may contain any one of the animals from that NFT Drop, including potentially the rarest.

To breed two adult NFT Animals, the user selects each animal in the user interface within the dApp, then initiates the process by ``feeding'' the designated amount of \$ZOO tokens. This mints a random NFT Egg, which the user can sell on the marketplace for profit, hold onto and sell later, transfer to any other user, or incubate and hatch into the random NFT Animal within.
